This thesis concerns the design and evaluation of wireless communication performance in real-time embedded systems using the ESP32 microcontroller. As part of the work, an original research platform named ESPy-Nexus was developed in a Hardware-in-the-Loop architecture, combining a real embedded device with a computer-based test engine, a database, and a web analytical interface. The main objective was to empirically determine how communication protocol, network topology, payload size, and transmission frequency affect reliability, timeliness, and throughput. The study includes automated experiments covering SERIAL, UDP, TCP, and WebSocket protocols across multiple network configurations. The collected data were processed into QoS metrics such as PDR, jitter, goodput, blackout time, relative delay, and clock drift. The results show that no single protocol is universally optimal: UDP provides better responsiveness and higher throughput at the cost of packet losses, while TCP and WebSocket better preserve stream integrity but may introduce significant latency under heavy load. The work provides a practical foundation for future research on the selection of communication protocols for real-time control and telemetry systems.

Keywords: Python; ESP32; wireless communication; QoS; UDP; TCP; WebSocket; Hardware-in-the-Loop; Wi-Fi networks; real-time systems; telemetry; network performance
